\documentclass[11pt,letterpaper]{article}
\usepackage[T1]{fontenc}
\usepackage{lmodern,amsmath,amssymb,amsthm,mathtools}
\usepackage[margin=0.92in,headheight=15pt]{geometry}
\usepackage{microtype,booktabs,longtable,tabularx,array}
\usepackage{xcolor,enumitem,fancyhdr,xurl,seqsplit}
\usepackage[colorlinks=true,linkcolor=blue!45!black,urlcolor=blue!50!black,
citecolor=blue!45!black]{hyperref}
\setlength{\parindent}{0pt}
\setlength{\parskip}{5pt plus 1pt}
\setlist{itemsep=2pt,topsep=3pt}
\pagestyle{fancy}
\fancyhf{}
\fancyhead[L]{\small Compatible decoder reductions for Con(PA)}
\fancyhead[R]{\small 29 October 2026}
\fancyfoot[C]{\thepage}
\newtheorem{theorem}{Theorem}[section]
\newtheorem{lemma}[theorem]{Lemma}
\newtheorem{proposition}[theorem]{Proposition}
\newcommand{\PA}{\mathrm{PA}}
\newcommand{\Con}{\operatorname{Con}}
\newcommand{\Halt}{\operatorname{Halt}}
\newcommand{\Sub}{\operatorname{Sub}}
\newcommand{\DU}{\mathcal D_U}
\newcommand{\row}[3]{#1\mathrm{#2}#3}
\newcommand{\file}[1]{\path{#1}}
\newcommand{\hash}[1]{\texttt{\seqsplit{#1}}}
\hypersetup{pdftitle={A 209-state candidate for Con(PA): compatible contextual decoder reductions},
pdfauthor={Research prepared for Vladimir Reshetnikov},
pdfsubject={Comparison of three U211 proposals and explicit U210 and U209 constructions}}
\title{\vspace{-1cm}\textbf{A 209-state candidate for the consistency\\of Peano arithmetic}\\[7pt]
\large Making three contextual decoder reductions compatible}
\author{Research prepared for Vladimir Reshetnikov}
\date{29 October 2026}

\begin{document}
\maketitle
\thispagestyle{empty}
\begin{abstract}
We compare the three 211-state proposals in the specified ProveIt incoming
directory and construct a binary, single-tape machine with
\textbf{209 working states and 418 transitions}, including its blank-tape
initializer. A separately verified 210-state intermediate is also supplied.
The three input machines have the same numerical proof-search program, but
their rooted labeled transition graphs are pairwise nonisomorphic. Their
savings remove different native states and cannot be stacked by copying
patches without preserving their entry conditions.

The new construction first combines the Research and Rotated savings by eight
closed source reassociations. It then changes a 720-primitive region across
a decoder shared with the runtime, preserving the specific branches used by
that runtime entry. An eight-factor reassociation and three further rotations
make the Contextual saving compatible with the other two. Native states
181, 218, and 180 are all absent from the final table.

Fresh symbolic derivations establish all 5,084 guarded primitive contracts
for the final table, with arbitrary natural registers, arbitrary extra header
gap, arbitrary nonempty finite spare lists, and complete-tape endpoints.
Every literal row is covered. Independent finite checkers verify the source
transformation, entry repair, local folds, and exact state counts; the complete
48,140-step blank initialization is unchanged. Composing these new arguments
with the retained numerical-calculus and unary-theory calibration yields the
ordinary mathematical, PA-formalizable endpoint
\(\PA\vdash\neg\Halt(U_{209},\mathrm{blank})\leftrightarrow\Con(\PA)\).
This package supplies mathematical proofs and executable certificates,
not a newly exported Lean, Rocq, MM0, or PA proof object for that endpoint.
\end{abstract}

\begin{center}
\fbox{\begin{minipage}{0.92\linewidth}
\textbf{Primary artifact:} \file{machine/u209.tm}. Working states \(0,\ldots,208\);
separate halt label \(209\); start state and head coordinate \(0\); all-zero
initial tape; one two-way infinite tape; alphabet \(\{0,1\}\); only left and
right moves. All \(209\cdot2=418\) working-state/symbol pairs are defined.

\textbf{Secondary artifact:} \file{machine/u210.tm}. The same conventions,
with working states \(0,\ldots,209\) and halt label \(210\).
\end{minipage}}
\end{center}

\section{Result, scope, and exact identities}

The requested threshold is surpassed by two states. Relative to the common
212-state parent, all three distinct published one-state savings now coexist:
the Research deletion of 181, the Rotated deletion of 218, and a contextual
deletion of 180 using a newly arranged keeper 216. No extra alphabet symbol,
prewritten input, additional tape, or omitted initializer is used.

A literal row \(\texttt{q b w D r}\) reads \(b\) in state \(q\), writes \(w\),
moves one cell in direction \(D\), and enters \(r\). Entering the separate
halt label ends the computation after the write and move. Native tables keep
sparse inherited labels and halt label 254; these labels are names, not counts.
The dense tables are exact checked renamings. Every working state is reachable
in the transition graph. This does not assert that every state is visited on
the particular blank run.

\begin{table}[ht]
\centering
\begin{tabular}{lrr}
\toprule
Property & \(U_{210}\) & \(U_{209}\)\\
\midrule
Working states & 210 & \textbf{209}\\
Transition rows & 420 & \textbf{418}\\
Source primitive occurrences & 3,909 & 3,909\\
Distinct syntactic subtrees & 197 & 197\\
Maximum source-address depth & 41 & 41\\
Maximum consecutive zeros in an address & 6 & 6\\
Nonhalt source successor edges & 5,083 & 5,083\\
Guarded primitive cases, including halt & 5,084 & 5,084\\
\bottomrule
\end{tabular}
\caption{The savings come from contextual control-state sharing, with the same
source program and the same number of distinct syntactic subtrees.}
\end{table}

The primary dense-table SHA-256 is
\begin{center}\small
\hash{bda91bd40880f96edef104246d22797fe281355cf03610ab7e93e7923be8aa16}.
\end{center}
Its native-table and source-map hashes are, respectively,
\begin{center}\small
\hash{4564c4876f9b18725bfb3505ec1f50d1ebce572502f2c6381ad673ecfdc3c3e1},\\
\hash{912724f29b0de28b3326c710fdf1f3863ebd07058c6db1a9b4bbaffb3ae5e3f3}.
\end{center}
The complete 209-state table appears in Appendix~\ref{app:table}.
Both state bijections, both literal tables, and all intermediate tables needed
by the local proofs are included as machine-readable files.

\begin{theorem}[New halting-preserving construction]\label{thm:new}
Let \(A\) be the exact Research U211 machine retained in this package.
The source identities, final-table primitive simulations, and common
blank initialization described below give PA-formalizable proofs of
\[
 \Halt(U_{209},\mathrm{blank})
 \ \longleftrightarrow\
 \Halt(U_{210},\mathrm{blank})
 \ \longleftrightarrow\
 \Halt(A,\mathrm{blank}).
\]
\end{theorem}

\begin{theorem}[Consistency endpoint with the retained calibration]\label{thm:endpoint}
Composing Theorem~\ref{thm:new} with the retained numerical proof-search and
Theory U calibration gives, in ordinary mathematics,
\[
 \PA\vdash
 \neg\Halt(U_{209},\mathrm{blank})
 \ \longleftrightarrow\ \Con(U)
 \ \longleftrightarrow\ \Con(\PA).
\]
\end{theorem}

Here \emph{retained calibration} names mathematical dependencies rather than a
new object-theory assumption. In particular the biconditional does not assume
\(\Con(\PA)\). Actual unconditional nonhalting is not established.
The distinction between an ordinary PA-formalizable argument and an exported,
kernel-checked PA endpoint is maintained throughout this report.

\section{Comparison of the three supplied proposals}\label{sec:comparison}

\subsection{Snapshot and literal comparison}
The comparison is pinned to ProveIt commit
\begin{center}\small
\hash{da4ab9dd68ccfc47ac852118601eaf1fdae53795}.
\end{center}
Its \file{docs/incoming/} directory contained the three U211 archives and
\file{pa_unary_power_cover.zip} used here~\cite{snapshot}.
All four original archives are retained unchanged in \file{inputs/}.
Their SHA-256 values, original paths, and sizes are recorded in
\file{provenance.json}. The table/source snapshots used by the new construction
are checked directly against the corresponding ZIP members.

\begin{table}[ht]
\centering\small
\begin{tabularx}{\linewidth}{@{}lrrX@{}}
\toprule
Proposal & Rotations & Deleted label & Principal mechanism\\
\midrule
Contextual & 6 & 180 &
Predecessor-tag fold into 163 through helper 211;
also repairs the backward entry \(171/0\) into state 7.\\
Research & 9 & 181 &
Predecessor-tag fold into 220 through helper 129;
retains the earlier helper 211.\\
Rotated & 11 & 218 &
Makes every entry to 218 write a zero immediately left of the head,
then replaces 218 by the existing state 157.\\
\bottomrule
\end{tabularx}
\caption{Distinct 211-state constructions from the same U212 parent.}
\end{table}

All three contain the same ordered 3,909 primitive operations and operands.
Both source-control resolvers produce the same complete ordinal control-flow
graph. However, their native tables differ, and canonical breadth-first
numbering from state 0 shows that their rooted transition graphs are pairwise
nonisomorphic when read symbols, write symbols, and directions are preserved.
These are different machines, not three filenames for one table.

\begin{table}[ht]
\centering\small
\begin{tabular}{lrr}
\toprule
Compared pair & Different common native rows & Different source addresses\\
\midrule
Contextual / Research & 28 & 1,275\\
Contextual / Rotated & 31 & 3,096\\
Research / Rotated & 29 & 3,096\\
\bottomrule
\end{tabular}
\caption{The row counts exclude the two outgoing rows of each differently
deleted label. Full inventories are in \texttt{comparison/census.json}.}
\end{table}

The common parent has native SHA-256
\begin{center}\small
\hash{386ce16259d4c97e4af31b3ba494feaff11a0a1a066e0235dddc154b4cea38c4}.
\end{center}
The three U211 native identities are listed below in abbreviated form;
the full hashes are retained in the comparison files.
\[
 \begin{array}{ll}
 \text{Contextual:}&\texttt{ac75e2fa934e89b3\ldots}\\
 \text{Research:}&\texttt{7edae566b4c2e5c3\ldots}\\
 \text{Rotated:}&\texttt{72234980bd604524\ldots}.
 \end{array}
\]

\subsection{Verification quality and the selected parent}
Each original package was copied to an isolated working directory and its
finite replay command was rerun successfully: Contextual's default reproduction,
Research's \texttt{--quick} reproduction, and Rotated's default verifier.
Those fresh receipts and logs are included. The original full symbolic runtime
suites were not rerun for all three inputs; their packaged full-runtime
receipts remain prior evidence. In contrast, the complete final-table runtime
derivations for both new candidates were freshly regenerated.

The Research table is the direct parent of the new work because its tagged
fold can be preserved while the 218 predecessor condition is rearranged.
It also gives a particularly transparent 210-state intermediate without the
171-entry repair. The Contextual package contributes an independently useful
retagging and repair argument; Rotated contributes the exact 218-to-157
two-cell identity. None of these merits implies that its entire patch can be
copied into another source layout without rechecking the entry conditions.

The common consistency calibration is not a reason to treat all three as
already covered by K215's separate formal endpoint. The linked
K215/MM0 project proves a statement about a different literal table and
ZFC~\cite{k215pa}. Its result motivates the architecture but does not certify
the newly supplied PA machines.

\section{Theory U and the arithmetic program actually preserved}\label{sec:theory}

\subsection{The unary schema}
Theory U uses classical, equality-free first-order logic on nonempty domains,
with one binary relation \(x\in a\). Put
\[
 \Sub(x,a)\;:\!\Longleftrightarrow\;\forall z\,(z\in x\to z\in a).
\]
For every formula \(\varphi(x)\) with no other free variable, its sole schema is
\[
 \forall a\,\exists b\,
 \left[
   \forall x\,(x\in b\to\varphi(x))
   \ \land\
   \forall x\,(\Sub(x,a)\to(\varphi(x)\to x\in b))
 \right].
 \tag{U}
\]
Variables are freshened before insertion. The source \(a\) is not a permitted
free matrix parameter. This is a restriction on free roles, not a restriction
to formulas containing only one variable symbol.

The schema supplies a weak cover: all members of \(b\) satisfy the matrix, and
every subset of \(a\) satisfying it belongs to \(b\). It does not require every
member of \(b\) to be a subset of \(a\), and it does not postulate ambient
extensionality. Replacing the matrix by
\(\Sub(x,a)\land\varphi(x)\) would introduce the forbidden free source \(a\);
that apparent shortcut is not available in U.

\subsection{Why the retained U/PA proof is substantive}
The supplied unary manuscript~\cite{unary} contains a detailed ordinary
interpretation argument. Its upper direction interprets membership as binary
bit membership in PA. Bit inclusion implies \(x\le a\); for each fixed
arithmetical matrix, PA's induction constructs the finite characteristic code
on the interval through \(a\). This is formula-specific finite comprehension,
not a uniform truth predicate, and works inside arbitrary PA models.

For the lower direction, equality is interpreted by coextensionality.
The numerical representatives explicitly include saturation and suitable
first-order order/endpoint conditions. Local predecessor facts establish
induction before total successor existence is used. Comparison and least
failing numerical tuples then allow parameters to be bound inside a legal
unary matrix. Bounded numerical comprehension follows, with its support and
invariance premises stated explicitly. This provides successors, numerical
powers, pair and graph codes, and canonical finite histories for addition and
multiplication. The order of these steps avoids a circular appeal to total
successors.

The manuscript also describes uniform syntactic proof templates: freshening,
substitution, relativization, replacement of terms by functional graphs, and
concatenation of finite derivations. These are primitive recursive syntax
operations, with correctness justified by induction in PA. An independent
review here found no concrete defect in that dependency chain. The review
does not claim to have implemented a complete arbitrary-proof translator or
checked an exported endpoint in a proof assistant.

\subsection{The machine-facing calculus is a different presentation}
The preserved source program searches the inherited numerical calculus
\(\DU\). It uses truth code \(0\), falsity code \(1\), native Cantor-coded
constructors, seven logical rows, a dominated propositional row, modus ponens,
generalization, and a gated cover/vacuity bundle.
Its machine-facing theory allows additional matrix parameters and uses global
irreflexivity. It is calibrated by finite-refutation translations to the
unary U presentation, rather than having literally the same axiom strings.

This distinction is essential. The separate positive dyadic emitter in the
unary-theory ZIP is not the program implemented by these U211, U210, or U209
tables. Its operation counts cannot be used as a runtime or state-count
certificate for them.

The retained role-emitter proof handles small-code aliases, the truth/falsity
codes, the gate, and the possible coincidence of output and source labels.
In the coincident-label case, irreflexivity changes the relevant guard to
emptiness; an empty source supplies the required cover. Padding and fresh
labels supply completeness for arbitrary desired instances.
The new reductions leave this numerical calculus and every ordered source
primitive unchanged.

\section{Source equivalence and contextual decoding}\label{sec:method}
The source is a binary tree whose leaves carry \(\mathrm{CI}\) (increment),
\(\mathrm{CD}\) (guarded decrement), \(\mathrm{CB}\) (forward branch),
\(\mathrm{CC}\) (backward branch), or \(\mathrm{CH}\) (halt).
There are respectively 1,558, 1,175, 80, 1,095, and one such occurrences.
The leaf address encodes control context, so retaining the printed operation
sequence alone is insufficient.

For ordinal \(i\), let \(G(i)\) record the primitive, its operand, and its ordered
successor ordinals. A decrement has zero and positive successors; halt has
none. The first resolver computes successors by binary increment and suffix
cutting, followed by zero-padded leaf decoding. The second resolver uses the
ordinal intervals of enclosing subtrees. Both check all 5,083 nonhalt
successor edges. The arithmetic resolver additionally rejects stray one bits
after the decoded target leaf.

A factor is closed and nonhalting when its ancestor-relative branches are
resolved within the factor and it contains no halt. Closure is checked in
two ways: recursive boundary signatures and independent leafwise counts of
the relevant enclosing zero/one edges. Association of closed sequential
factors preserves the primitive word and resolved ordinal control graph.

\begin{lemma}[Source conjugacy]\label{lem:source}
If two complete source trees have the same ordered primitive word and the
same \(G(i)\) for every ordinal, then from a common ordinal and numerical
register store they have identical stores and ordinals at every source step.
They halt after the same number of source steps.
\end{lemma}
\begin{proof}
Induct on the source-step count. Increment changes the same register in both
programs. A decrement tests the same value, so both programs choose the same
branch and make the same register change. The two branch operations choose
the checked common successor; halt occurs in the same ordinal. No bound on
the register values is used.
\end{proof}

A pure physical decoder has rows \((0,R,a)\) and \((1,R,b)\).
A tagged physical decoder may implement different logical source subtrees
depending on a bit left by its incoming transition. Therefore a source rewrite
must preserve the relevant entry context. This is the key to the new result:
one may alter a branch of a shared decoder that a runtime call cannot execute,
while preserving all branches used by that runtime call.

The source theorem is separate from the final literal-machine theorem.
Some construction intermediates have seven consecutive zero address bits,
which collide with the inherited origin detector. Their source control graphs
are still equivalent, but those intermediate tables are not claimed to
implement a complete correct runtime. The final sources have maximum zero
run six and pass the complete runtime verification directly.

\section{First construction: Research U211 to U210}\label{sec:u210}

\subsection{Move the continuation while preserving the Research selector}
Starting with the Research table, apply these eight left associations:
\[
\begin{array}{rrrr}
(162,203,0),&(203,223,0),&(223,188,0),&(188,177,0),\\
(177,89,0),&(89,174,0),&(174,217,0),&(217,163,0).
\end{array}
\]
Here \((p,q,0)\) changes \(((A,B),C)\) at decoder \(p\), whose zero child is
\(q\), into \((A,(B,C))\). The builder verifies pure decoder rows, the complete
unique incoming edge to the inner state, and closure of each of the three
source factors. The independent checker recovers the affected occurrences
by literal source-prefix execution and checks both complete control graphs
after every source change.

These associations carry continuation 116 to the right of state 163.
The Research selector rows at 220 and 129 remain intact.
The old 181 state remains deleted. The final remaining entrances to 218 are
exactly
\[
 158/0:\row0R{218},\qquad
 163/0:\row0R{218},\qquad
 211/0:\row0R{218}.
 \tag{1}
\]
Consequently each visit to 218 from state 0 has a zero immediately left of
the head. This is a direct consequence of the preceding write/move, valid
for every initial tape, rather than a property guessed from a blank trace.

\subsection{The 218-to-157 local identity}
The old rows are \(218/0=\row0R{90}\) and \(218/1=\row1R{86}\).
Redirect all targets of 218 to 157 and delete 218. With the original head at
coordinate zero, the existing table executes
\[
\begin{array}{c|l}
\text{current bit}&\text{replacement state/head path}\\
\hline
0&157@0\to67@{-1}\to151@0\to90@1\\
1&157@0\to148@{-1}\to186@0\to86@1.
\end{array}
\tag{2}
\]
Only cells \(-1,0\) are read. The predecessor is zero by (1);
each write preserves its bit, and the exterior is untouched.

\begin{lemma}[U210 final local fold]
From every common initial tape in state 0, the 211-state prefold table and
the delivered U210 halt equivalently. Each old 218 step becomes three new
steps; every other old step becomes one new step.
\end{lemma}
\begin{proof}
Map 218 to 157 and fix all other labels. For ordinary old rows, write and move
are unchanged and the target is mapped by this projection. For old state 218,
use the exhaustive two cases (2) and the incoming-edge invariant (1).
Every replacement block has positive length and no internal halt.
Induction on old steps gives the simulation, including equality of the whole
tape at its block boundaries. Conversely, if the new machine halts after
\(t\) steps, simulating \(t\) old steps would pass that point unless the old
machine had already halted; positivity and the absence of internal halt
exclude a spurious new halt.
\end{proof}

For any halting old run, or aligned finite run prefixes, this final fold gives
\(T_{\rm new}=T_{\rm old}+2k\), where \(k\) counts old visits to 218, and
\(T_{\rm old}\le T_{\rm new}\le3T_{\rm old}\).
This clock bound is not asserted across the preceding address changes.

The exact U210 dense hash is
\begin{center}\small
\hash{9020e082ef6995d13cd571c325f4d83bf615f4d0a82e77c580353c1127a0d596}.
\end{center}
Its full final-table runtime and blank startup were checked independently of
the state-count calculation.

\section{Second construction: a compatible U209}\label{sec:u209}

\subsection{The shared helper that prevents a direct patch transplant}
In U210, state 129 is entered both by a right-moving source decoder and by
the backward call \(220/0=\row0L{129}\). Its two incoming rows are
\(215/1=\row1R{129}\) and \(220/0=\row0L{129}\).
Treating it as an ordinary singly entered source decoder would lose one
of these uses.

The runtime call begins after 220 has just read and preserved zero.
If the predecessor bit is \(t\), its path is
\[
 220@0\to129@{-1}
 \to\begin{cases}158@0,&t=0,\\204@0,&t=1,\end{cases}
 \to\begin{cases}157@1,&t=0,\\147@1,&t=1.\end{cases}
 \tag{3}
\]
Thus the call uses 129's two rows, but only the zero rows of 158 and 204.
The row \(204/1\) is available for a source-layout change without altering
this runtime block. This selective preservation is broader than requiring
the entire shared decoder to be frozen or to have a single entrance.

\subsection{A closed eight-factor reassociation}
At the U210 source prefix
\[
 \texttt{0011001110101011111111}
\]
state 174 represents primitive ordinals \([3016,3736)\), a region of length 720.
Write its eight closed factors in their unchanged order as follows.

\begin{table}[ht]
\centering\small
\begin{tabular}{cl r ll}
\toprule
Factor & Physical entry & Length & Old suffix & New suffix\\
\midrule
\(A\)&98&26&00&00\\
\(B\)&158&134&010&010\\
\(C\)&147&52&0110&0110\\
\(D\)&146&108&01110&0111\\
\(E\)&220, zero tag&136&011110&10\\
\(F\)&90&27&011111&1100\\
\(G\)&211&103&10&1101\\
\(H\)&163&134&11&111\\
\bottomrule
\end{tabular}
\caption{The zero tag in factor \(E\) is part of its physical decoding context.}
\end{table}

Replace the source shape by
\begin{align*}
 &((A,(B,(C,(D,(E,F))))),(G,H))\\
 &\hspace{2cm}\longmapsto
 ((A,(B,(C,D))),(E,((F,G),H))).
\end{align*}
Every factor is closed and nonhalting. The word and both complete ordinal
control graphs are preserved. Physically, this is implemented by seven row
changes:
\[
\begin{array}{c|c|c}
\text{row}&\text{before}&\text{after}\\
\hline
204/1&\row1R{214}&\row1R{146}\\
214/0&\row0R{146}&\row0R{216}\\
214/1&\row1R{216}&\row1R{163}\\
216/0&\row0R{220}&\row0R{90}\\
216/1&\row1R{90}&\row1R{211}\\
217/0&\row0R{211}&\row0R{220}\\
217/1&\row1R{163}&\row1R{214}
\end{array}
\tag{4}
\]
No state is added. State 216 now decodes the pair \((90,211)\).

The zero-tag entrance to 220 moves from \(216/0\) to \(217/0\); it still
writes zero and moves right. Its one-tag entrance \(185/1\) is unchanged.
All rows used by (3) are unchanged. In particular \(204/0\) is preserved
while \(204/1\) changes. The source entry and backward runtime entry are
therefore justified separately. Exhaustive final-prefix and primitive-runtime
checks additionally validate the actual literal table.

\subsection{Retag state 180 and repair the arithmetic entrance}
Apply the three right associations
\[
 (206,199,1),\qquad(205,206,1),\qquad(7,205,1),
 \tag{5}
\]
where side 1 changes \((A,(B,C))\) to \(((A,B),C)\).
The inner decoder has its specified unique incoming edge in each case, and
the three factors are closed. Both complete ordinal control graphs remain
unchanged.

State 7 also has the backward arithmetic entrance \(171/0=\row0L7\).
Repair this one row to
\[
                     171/0=\row0L{205}.
 \tag{6}
\]
For a current zero at coordinate zero, split on the predecessor:
\[
\begin{array}{c|l|l}
\text{left bit}&\text{old path}&\text{repaired path}\\
\hline
0&171@0\to7@{-1}\to86@0&
 171@0\to205@{-1}\to86@0\\
1&171@0\to7@{-1}\to205@0\to164@1&
 171@0\to205@{-1}\to206@0\to164@1.
\end{array}
\tag{7}
\]
Both paths preserve the two cells and read no exterior cell. This establishes
the entry repair on every such configuration.

The first two stages of (5) have maximum zero run seven; their intermediate
tables are not asserted to be complete runtimes. The final stage restores
maximum zero run six. The independent receipt records any failed intermediate
prefix check as \texttt{NOT\_A\_RUNTIME\_INTERMEDIATE}; it checks every prefix
of the final table successfully.

\subsection{Delete 180 using a five-step contextual path}
In the 210-state prefold table \(R\),
\[
\begin{array}{c|cc}
q&\delta(q,0)&\delta(q,1)\\
\hline
216&\row0R{90}&\row1R{211}\\
180&\row0R{152}&\row1R{211}.
\end{array}
\]
Its complete incoming inventories are
\[
 \operatorname{In}(216)=\{214/0:\row0R{216}\},\qquad
 \operatorname{In}(180)=
 \{132/1:\row1R{180},\,199/1:\row1R{180}\}.
 \tag{8}
\]
Neither state is initial. Therefore the old state is encoded by the bit
immediately left of the head: zero for 216 and one for 180.

Let \(\rho(180)=216\) and \(\rho(q)=q\) otherwise. Redirect both incoming targets
of 180 to 216, delete its two rows, and replace
\[
                           216/0\quad\text{by}\quad\row0L{211}.
 \tag{9}
\]
Keep \(216/1=\row1R{211}\). The resulting machine \(C\) has 209 states.

\begin{table}[ht]
\centering\small
\begin{tabularx}{\linewidth}{@{}rrrrX@{}}
\toprule
Old state & Left bit & Read & New steps & New state/head path\\
\midrule
216&0&0&5&\(216@0,211@{-1},157@0,67@{-1},151@0,90@1\)\\
180&1&0&3&\(216@0,211@{-1},184@0,152@1\)\\
216&0&1&1&\(216@0,211@1\)\\
180&1&1&1&\(216@0,211@1\)\\
\bottomrule
\end{tabularx}
\caption{Exhaustive whole-tape identities for the final fold. Each old block
has one step. Only the predecessor and current cell are read.}
\label{tab:fold}
\end{table}

The five-step case is the compatibility point: 218 has already been deleted,
so the zero predecessor branch must pass through its existing 157 detour.
Using the earlier three-step formula here would be wrong.

\begin{lemma}[Final fold, arbitrary exterior]\label{lem:209fold}
From state 0 on every common initial tape, \(R\) and \(C=U_{209}\) halt
equivalently. For a halting old run,
\[
 T_C=T_R+4n_0+2n_1,\qquad T_R\le T_C\le5T_R,
\]
where \(n_0\) counts old read-zero visits to 216 and \(n_1\) counts old
read-zero visits to 180.
\end{lemma}
\begin{proof}
Inventory (8) establishes the tag invariant whenever either old state is
entered. In the four exhaustive cases of Table~\ref{tab:fold}, direct literal
execution leaves both displayed cells unchanged, reads no other cell, and
ends with the same head and the projected old target. No internal state is
the halt label. Every remaining old row commutes with \(\rho\), preserving
its write and direction. Induction gives a positive block simulation.
A new finite halt cannot occur inside a block, and each block consumes at
least one new step, so the converse follows by the finite-horizon argument
used for U210. Summing the four block lengths gives the clock identity.
For arithmetic codes of finite-support tapes these are finite induction
arguments formalizable in PA.
\end{proof}

The all-initial-tape assertion and time bound concern \(R\to C\) alone.
They are not extended to the preceding source-address transformations.

Relative to U210, U209 changes 15 retained rows and deletes both rows of 180.
It changes 847 source addresses, including 508 changed by (4), while retaining
the same 197 distinct syntactic subtrees.

\section{Universal runtime and exact initialization}\label{sec:verification}

\subsection{What the symbolic contracts quantify over}
A canonical frame at native dispatch state 84 contains a left marker, a
zero-padded source address, a header at distance \(56+g\), nine unary-coded
primary registers, and a nonempty finite spare-register list, with zero
exteriors. The physical register permutation is
\[
                         (6,2,0,4,3,7,8,5,1).
\]
A register value \(v\) is represented by the block
\(B(v)=0\,1^{v+1}\). The symbolic inputs range over arbitrary natural values
of all nine logical registers, an arbitrary natural extra gap \(g\), an
arbitrary first spare value, and an arbitrary remaining finite spare list.

For each source site, the verifier derives the exact primitive effect and
successor frame. Conditional decrement is split into zero and positive
cases, with the positive value written \(v+1\).
Increment appends the specified spare block, decrement and source arithmetic
adjust the header/address representation as required, and branch operations
leave the logical register store fixed. Nonhalt targets compare the complete
oriented tape terms, head displacement, and dispatch state. Halt requires
entry into native label 254.

The distinction between binary successor words and leaf addresses is checked:
a target word may contain additional zero padding after its first decoded
leaf, but it must contain no stray one bits there. An independent calculation
confirmed agreement with ordinal source successors in all 5,083 nonhalt
cases for each final candidate.

\subsection{The proof rules used by the checker}
The inherited engine is copied byte-for-byte from the Rotated package and
its provenance is retained. It uses three execution rules:
\begin{enumerate}
\item A literal step applies a row of the actual table and updates the
represented tape, state, motion, and clock.
\item A unary scan accelerates a literal self-loop in the scan direction
across \(b^n\), for arbitrary natural \(n\). Induction on \(n\) justifies
the resulting block and its exact duration.
\item A finite-list traversal uses a one-block theorem first rederived from
the actual table with opaque outside sentinels. Each such theorem returns
to the same state and direction, handles an arbitrary block value \(v\),
takes \(v+2\) steps, and preserves the exterior. Induction then traverses an
arbitrary finite list.
\end{enumerate}
The four one-block lemmas are regenerated on both final tables.
Normalization uses literal word identities and the valid finite-list identity
\(C(L)1=1B(L)\), where \(C(v)=1\,0\,1^v\), together with its reversal.
Its base case and one-block algebra are checked; list induction supplies
the arbitrary-list form.

Clock and motion expressions use thirteen affine slots: a constant, the gap,
nine primary values, the first spare value, and the total width of the
remaining spare list. Every accepted primitive duration has nonnegative
coefficients and strictly positive constant. These are symbolic induction
certificates, not numerical tests at 5,084 sampled register settings.
The engine is ordinary Python, so the report makes its small set of
mathematical rules and their trust boundary explicit rather than calling it
a proof-assistant kernel.

\begin{table}[ht]
\centering
\begin{tabular}{lrr}
\toprule
Fresh final-table verification & U210 & U209\\
\midrule
Guarded contracts & 5,084 & 5,084\\
Symbolic inference applications & 1,249,857 & 1,248,748\\
Literal rows covered & 420/420 & 418/418\\
Minimum duration constant & 5 & 5\\
Final internal-prefix checks & 3,908 & 3,908\\
Maximum steps in those prefix checks & 49 & 49\\
Semantic corruptions rejected & 9 & 11\\
\bottomrule
\end{tabular}
\caption{Both complete runtime suites were freshly executed.
Symbolic inference counts are certificate-replay counts, not a claim about
the actual proof search's running time.}
\end{table}

\subsection{The exact blank startup}
Both new tables have the same complete 48,140-transition blank prefix as
the Research U211 parent. No changed row is visited. It ends at
\[
 q=84,\qquad h=-2,\qquad \text{header}=54,\qquad
 \text{source ordinal}=136,
\]
with source address \(\texttt{001100110101010}\), logical register store
\[
                         (1,0,13,0,0,1,0,0,1),
\]
and 43 spare registers of value one. The checker compares the entire set
of 120 nonzero tape cells, so the assertion includes the zero exterior,
not just a finite viewing window.

The same ordinal/store is reached after 141 source primitives from ordinal
106 with initial store
\[ (1,0,4,0,0,1,0,0,1). \]
This finite source bridge is replayed.
The inherited startup-suffix proof supplies the further logical meaning:
the incoming cache value 1 is used arithmetically and is not assumed to be
a theorem; the suffix creates nine theorem-valued heads above opaque tail 4.

That tail would be unsafe if exposed. The retained scheduler argument bounds
every command-prefix debt by two, leaving a permanent margin of seven
explicit heads. A ghost explicit stack length handles the pairing alias
\(P(0,0)=0\); a numerical stack-length decoder is not assumed. The retained
strong proof compiler works above arbitrary underlying stacks and produces
cofinally many successful candidate codes. These arguments provide soundness
and completeness after the actual startup and finite search cutoff.
Neither is replaced by the finite 48,140-step test.

\subsection{Independent checks and negative controls}
The finite checkers do not import the construction builder. They derive all
3,908 internal source-prefix bindings from literal execution, use two
complete CFG resolvers, test closure leaf by leaf, check full incoming
inventories and every table row, and verify dense renaming.
The U209 audit also checks the two arithmetic entry identities and four
Research-selector cases on the final table.

The nine U210 and eleven U209 negative controls deliberately damage writes,
moves, helper targets, endpoint rows, tag inventories, or the 171 repair.
They are rejected by semantic checks rather than only by a hash mismatch.
Both independent checkers produce byte-identical receipts normally and under
\texttt{python3 -O}, because their obligations use explicit exceptions.
The full symbolic engine uses assertions and rejects optimized Python.

\section{Proof of the machine-to-consistency composition}
Lemma~\ref{lem:source} identifies all source executions by their ordinals and
logical stores. The final runtime contracts show that every source primitive
is simulated by a finite nonempty block of the corresponding literal machine.
Their complete target frames maintain the induction invariant, including the
gap and spare-register list. The unique halt primitive reaches the separate
halt label; the other contracts return to dispatch with the prescribed
successor and exclude earlier termination.

The exact startup places both final machines in the calibrated common
interior source configuration. Consequently the source simulation applies
to their actual blank runs. The inherited parent analysis gives the same
source interpretation for the selected Research U211. This proves
Theorem~\ref{thm:new}. It is not necessary to claim that every temporary
reassociation table implements its temporary source.

The source proof search is unchanged, so its retained calibration gives
\[
 \Halt(U_{209},\mathrm{blank})
 \ \longleftrightarrow\ (\DU\vdash1)
 \ \longleftrightarrow\ (\PA\vdash\bot)
\]
as the corresponding arithmetized equivalence. The source compiler,
protected-tail proof, and finite-refutation translations discussed in
Section~\ref{sec:theory} are the explicit inherited links.
Negating the existence of a finite refutation yields
Theorem~\ref{thm:endpoint}.

Finite tables, bounded local paths, syntax transformations, and natural/list
induction are all arithmetizable in PA. This explains the claimed
formalizability. It does not turn a Python replay into a checked exported
PA derivation. The archive contains no newly generated end-to-end proof object
for U209 in Lean, Rocq, or MM0.

\section{Reproduction, deliverables, and remaining work}
The package is self-contained and uses Python 3.10 or later with its standard
library. From its extracted root:
\begin{verbatim}
python3 reproduce.py --quick
python3 reproduce.py
python3 reproduce.py --both
\end{verbatim}
The quick mode checks all frozen hashes, reconstructs both machines, binds
the three comparison snapshots to the original archives, replays the two
independent finite audits and both startups, and checks the delivered runtime
records as data. It explicitly skips new universal derivations.
The default additionally regenerates all U209 contracts.
The \texttt{--both} mode regenerates both complete suites and requires exact
equality with the delivered per-case symbolic records. Runtime depends on
the computer; the fresh derivations here took a few minutes per table.

Generated replay outputs go under \file{build/}; delivered evidence is not
overwritten. The complete payload is hash-indexed in \file{MANIFEST.json}.
To rebuild this report, use a standard LaTeX installation:
\begin{verbatim}
pdflatex -interaction=nonstopmode -halt-on-error report.tex
pdflatex -interaction=nonstopmode -halt-on-error report.tex
\end{verbatim}
Verify the delivered manifest before rebuilding the PDF, whose metadata
can change with the TeX installation.

\begin{table}[ht]
\centering\small
\begin{tabularx}{\linewidth}{@{}lX@{}}
\toprule
Location & Contents\\
\midrule
\file{machine/} & Both dense and native tables and exact state bijections.\\
\file{source/}, \file{intermediates/} & Final source maps and immediate prefold tables.\\
\file{code/} & Deterministic builder, independent checkers, comparison, startup,
and the unchanged inherited symbolic engine.\\
\file{evidence/} & Fresh complete runtime cases and summaries, independent finite
receipts, startup frames, and source-stage diagnostics.\\
\file{comparison/} & Three exact input snapshots, pairwise census, and fresh
original finite replay logs.\\
\file{inputs/} & All four original archives and the short construction recipes.\\
\file{proofs/} & Detailed theory, runtime, comparison, and new-reduction audit notes.\\
\bottomrule
\end{tabularx}
\end{table}

The research contribution established here is relative to the supplied
proposals: previously separate savings have been made compatible through
source reorganization that respects the exact context of a shared helper.
No global minimality, independent record ratification, or priority claim
is made. The natural next verification milestone is to export the local
context certificates and source/runtime simulation to a small proof-assistant
kernel, and then connect that development to the retained arithmetic
calibration. The concrete machine and its current proof boundary are fully
specified without that additional work.


\clearpage
\appendix
\section{Complete dense transition table for U209}\label{app:table}
Start in state 0 on an all-zero tape. States \(0,\ldots,208\) are working
states; 209 is the separate halt label. An entry \(wDr\) writes \(w\), moves
in direction \(D\), and enters \(r\). Each displayed row supplies both
binary-symbol transitions. The equivalent plain-text table is
\file{machine/u209.tm}.
\begingroup
\setlength{\tabcolsep}{3.4pt}
\renewcommand{\arraystretch}{1.03}
\fontsize{8.3}{9.7}\selectfont
\begin{longtable}{@{}rll@{\quad}rll@{\quad}rll@{\quad}rll@{}}
\toprule
$q$ & $\delta(q,0)$ & $\delta(q,1)$ & $q$ & $\delta(q,0)$ & $\delta(q,1)$ & $q$ & $\delta(q,0)$ & $\delta(q,1)$ & $q$ & $\delta(q,0)$ & $\delta(q,1)$\\
\midrule\endhead
0 & $1\mathrm{R}92$ & $1\mathrm{R}0$ & 53 & $0\mathrm{L}40$ & $1\mathrm{L}53$ & 106 & $0\mathrm{R}61$ & $1\mathrm{R}42$ & 159 & $0\mathrm{R}19$ & $1\mathrm{R}171$\\
1 & $0\mathrm{L}1$ & $1\mathrm{L}2$ & 54 & $0\mathrm{R}22$ & $1\mathrm{R}54$ & 107 & $0\mathrm{R}121$ & $1\mathrm{R}8$ & 160 & $0\mathrm{R}164$ & $1\mathrm{R}187$\\
2 & $0\mathrm{L}3$ & $1\mathrm{L}2$ & 55 & $0\mathrm{L}11$ & $0\mathrm{R}50$ & 108 & $0\mathrm{R}101$ & $1\mathrm{R}122$ & 161 & $0\mathrm{R}159$ & $1\mathrm{R}96$\\
3 & $0\mathrm{L}5$ & $1\mathrm{L}2$ & 56 & $0\mathrm{R}60$ & $0\mathrm{L}16$ & 109 & $0\mathrm{L}208$ & $0\mathrm{L}95$ & 162 & $0\mathrm{R}125$ & $1\mathrm{R}164$\\
4 & $1\mathrm{R}6$ & $0\mathrm{L}4$ & 57 & $0\mathrm{R}69$ & $1\mathrm{R}70$ & 110 & $0\mathrm{R}124$ & $1\mathrm{R}104$ & 163 & $0\mathrm{R}97$ & $1\mathrm{R}162$\\
5 & $0\mathrm{L}9$ & $1\mathrm{L}2$ & 58 & $0\mathrm{R}61$ & $1\mathrm{R}71$ & 111 & $0\mathrm{R}115$ & $1\mathrm{L}72$ & 164 & $0\mathrm{L}190$ & $1\mathrm{R}142$\\
6 & $0\mathrm{R}6$ & $0\mathrm{R}10$ & 59 & $0\mathrm{R}54$ & $1\mathrm{R}59$ & 112 & $0\mathrm{R}165$ & $1\mathrm{R}181$ & 165 & $0\mathrm{L}66$ & $1\mathrm{R}173$\\
7 & $0\mathrm{R}190$ & $1\mathrm{R}159$ & 60 & $0\mathrm{R}49$ & $0\mathrm{R}25$ & 113 & $0\mathrm{R}160$ & $1\mathrm{R}184$ & 166 & $0\mathrm{R}146$ & $1\mathrm{R}172$\\
8 & $0\mathrm{R}14$ & $1\mathrm{R}13$ & 61 & $0\mathrm{L}28$ & $1\mathrm{L}72$ & 114 & $0\mathrm{R}148$ & $1\mathrm{R}146$ & 167 & $0\mathrm{R}199$ & $1\mathrm{R}201$\\
9 & $0\mathrm{L}15$ & $1\mathrm{L}2$ & 62 & $0\mathrm{L}117$ & $1\mathrm{L}30$ & 115 & $0\mathrm{R}11$ & $1\mathrm{R}129$ & 168 & $0\mathrm{R}161$ & $1\mathrm{R}143$\\
10 & $0\mathrm{L}16$ & $0\mathrm{R}10$ & 63 & $0\mathrm{R}73$ & $1\mathrm{R}63$ & 116 & $0\mathrm{L}131$ & $0\mathrm{L}1$ & 169 & $0\mathrm{R}148$ & $1\mathrm{R}87$\\
11 & $0\mathrm{L}26$ & $1\mathrm{R}128$ & 64 & $0\mathrm{R}97$ & $1\mathrm{R}150$ & 117 & $0\mathrm{R}39$ & $1\mathrm{R}72$ & 170 & $0\mathrm{R}161$ & $1\mathrm{R}96$\\
12 & $0\mathrm{R}19$ & $1\mathrm{R}20$ & 65 & $1\mathrm{R}185$ & $1\mathrm{L}53$ & 118 & $0\mathrm{R}61$ & $1\mathrm{R}132$ & 171 & $0\mathrm{R}181$ & $1\mathrm{R}122$\\
13 & $0\mathrm{L}89$ & $1\mathrm{L}31$ & 66 & $0\mathrm{R}144$ & $1\mathrm{R}145$ & 119 & $0\mathrm{R}61$ & $1\mathrm{R}133$ & 172 & $0\mathrm{L}193$ & $1\mathrm{R}60$\\
14 & $0\mathrm{L}11$ & $0\mathrm{R}22$ & 67 & $0\mathrm{R}76$ & $1\mathrm{R}209$ & 120 & $0\mathrm{R}70$ & $0\mathrm{L}16$ & 173 & $0\mathrm{R}145$ & $1\mathrm{R}88$\\
15 & $0\mathrm{L}23$ & $1\mathrm{L}2$ & 68 & $0\mathrm{L}197$ & $1\mathrm{R}31$ & 121 & $0\mathrm{R}14$ & $1\mathrm{R}132$ & 174 & $0\mathrm{R}152$ & $1\mathrm{R}202$\\
16 & $1\mathrm{L}1$ & $0\mathrm{L}1$ & 69 & $0\mathrm{R}41$ & $1\mathrm{R}62$ & 122 & $0\mathrm{R}21$ & $1\mathrm{R}134$ & 175 & $0\mathrm{R}140$ & $1\mathrm{R}84$\\
17 & $0\mathrm{R}24$ & $0\mathrm{R}32$ & 70 & $0\mathrm{L}89$ & $1\mathrm{R}78$ & 123 & $0\mathrm{R}151$ & $1\mathrm{R}189$ & 176 & $0\mathrm{R}175$ & $1\mathrm{R}75$\\
18 & $0\mathrm{R}98$ & $1\mathrm{R}113$ & 71 & $0\mathrm{R}55$ & $1\mathrm{L}79$ & 124 & $0\mathrm{R}137$ & $1\mathrm{L}20$ & 177 & $0\mathrm{R}126$ & $1\mathrm{R}169$\\
19 & $0\mathrm{R}28$ & $1\mathrm{R}29$ & 72 & $0\mathrm{R}80$ & $1\mathrm{R}81$ & 125 & $0\mathrm{R}130$ & $1\mathrm{R}180$ & 178 & $0\mathrm{R}138$ & $1\mathrm{R}116$\\
20 & $0\mathrm{R}30$ & $1\mathrm{R}31$ & 73 & $1\mathrm{R}0$ & $1\mathrm{R}63$ & 126 & $0\mathrm{R}204$ & $1\mathrm{R}200$ & 179 & $0\mathrm{R}97$ & $1\mathrm{R}51$\\
21 & $0\mathrm{R}24$ & $0\mathrm{R}59$ & 74 & $0\mathrm{R}74$ & $1\mathrm{R}82$ & 127 & $0\mathrm{R}75$ & $1\mathrm{R}176$ & 180 & $0\mathrm{R}168$ & $1\mathrm{R}51$\\
22 & $0\mathrm{R}33$ & $1\mathrm{R}22$ & 75 & $0\mathrm{L}145$ & $1\mathrm{L}141$ & 128 & $0\mathrm{R}75$ & $1\mathrm{R}206$ & 181 & $0\mathrm{R}17$ & $1\mathrm{L}131$\\
23 & $0\mathrm{L}34$ & $1\mathrm{L}2$ & 76 & $0\mathrm{R}85$ & $1\mathrm{R}86$ & 129 & $0\mathrm{R}203$ & $1\mathrm{R}157$ & 182 & $0\mathrm{R}143$ & $1\mathrm{R}149$\\
24 & $0\mathrm{L}35$ & $0\mathrm{R}0$ & 77 & $0\mathrm{R}20$ & $1\mathrm{R}38$ & 130 & $0\mathrm{R}12$ & $1\mathrm{R}141$ & 183 & $0\mathrm{R}76$ & $1\mathrm{R}116$\\
25 & $0\mathrm{R}36$ & $1\mathrm{R}25$ & 78 & $0\mathrm{R}89$ & $1\mathrm{L}132$ & 131 & $0\mathrm{R}30$ & $1\mathrm{R}69$ & 184 & $0\mathrm{R}106$ & $1\mathrm{R}124$\\
26 & $0\mathrm{R}37$ & $1\mathrm{R}38$ & 79 & $0\mathrm{L}108$ & $1\mathrm{R}56$ & 132 & $0\mathrm{R}39$ & $1\mathrm{L}124$ & 185 & $1\mathrm{L}40$ & $1\mathrm{L}109$\\
27 & $0\mathrm{R}39$ & $0\mathrm{L}40$ & 80 & $0\mathrm{R}49$ & $0\mathrm{R}33$ & 133 & $0\mathrm{R}153$ & $0\mathrm{L}16$ & 186 & $0\mathrm{R}175$ & $1\mathrm{R}200$\\
28 & $0\mathrm{R}41$ & $1\mathrm{R}42$ & 81 & $0\mathrm{R}91$ & $0\mathrm{L}42$ & 134 & $0\mathrm{R}154$ & $0\mathrm{L}16$ & 187 & $0\mathrm{R}182$ & $1\mathrm{R}183$\\
29 & $0\mathrm{R}43$ & $1\mathrm{R}44$ & 82 & $0\mathrm{R}178$ & $1\mathrm{R}77$ & 135 & $0\mathrm{R}88$ & $1\mathrm{R}136$ & 188 & $0\mathrm{R}66$ & $1\mathrm{R}205$\\
30 & $0\mathrm{R}45$ & $1\mathrm{R}46$ & 83 & $0\mathrm{R}96$ & $1\mathrm{R}12$ & 136 & $0\mathrm{R}197$ & $1\mathrm{R}140$ & 189 & $0\mathrm{R}140$ & $1\mathrm{R}139$\\
31 & $0\mathrm{R}21$ & $1\mathrm{R}47$ & 84 & $0\mathrm{R}88$ & $1\mathrm{R}97$ & 137 & $0\mathrm{R}24$ & $0\mathrm{R}48$ & 190 & $0\mathrm{R}84$ & $1\mathrm{R}191$\\
32 & $0\mathrm{R}48$ & $1\mathrm{R}32$ & 85 & $0\mathrm{R}207$ & $1\mathrm{R}21$ & 138 & $0\mathrm{R}99$ & $1\mathrm{R}37$ & 191 & $0\mathrm{R}157$ & $1\mathrm{R}186$\\
33 & $0\mathrm{R}32$ & $1\mathrm{R}33$ & 86 & $0\mathrm{R}116$ & $1\mathrm{R}100$ & 139 & $0\mathrm{R}172$ & $1\mathrm{R}148$ & 192 & $0\mathrm{R}64$ & $1\mathrm{R}7$\\
34 & $0\mathrm{L}74$ & $1\mathrm{L}2$ & 87 & $0\mathrm{R}174$ & $1\mathrm{R}167$ & 140 & $0\mathrm{R}96$ & $1\mathrm{R}159$ & 193 & $0\mathrm{R}150$ & $1\mathrm{R}127$\\
35 & $1\mathrm{R}35$ & $1\mathrm{L}1$ & 88 & $0\mathrm{R}102$ & $1\mathrm{R}103$ & 141 & $0\mathrm{R}175$ & $1\mathrm{R}166$ & 194 & $0\mathrm{R}18$ & $1\mathrm{R}94$\\
36 & $0\mathrm{R}50$ & $1\mathrm{R}36$ & 89 & $0\mathrm{R}80$ & $1\mathrm{R}79$ & 142 & $0\mathrm{R}12$ & $1\mathrm{R}159$ & 195 & $0\mathrm{R}150$ & $1\mathrm{R}173$\\
37 & $0\mathrm{R}17$ & $1\mathrm{R}111$ & 90 & $0\mathrm{R}101$ & $1\mathrm{R}196$ & 143 & $0\mathrm{R}88$ & $1\mathrm{L}166$ & 196 & $0\mathrm{R}104$ & $1\mathrm{R}96$\\
38 & $0\mathrm{R}13$ & $1\mathrm{R}52$ & 91 & $0\mathrm{R}105$ & $1\mathrm{R}106$ & 144 & $0\mathrm{R}88$ & $1\mathrm{R}158$ & 197 & $0\mathrm{R}90$ & $1\mathrm{R}88$\\
39 & $0\mathrm{L}11$ & $0\mathrm{R}36$ & 92 & $0\mathrm{R}65$ & $0\mathrm{R}0$ & 145 & $0\mathrm{L}152$ & $1\mathrm{L}163$ & 198 & $0\mathrm{R}200$ & $1\mathrm{R}156$\\
40 & $0\mathrm{L}4$ & $1\mathrm{L}53$ & 93 & $0\mathrm{R}27$ & $1\mathrm{R}107$ & 146 & $0\mathrm{R}84$ & $1\mathrm{R}88$ & 199 & $0\mathrm{R}96$ & $1\mathrm{R}123$\\
41 & $0\mathrm{R}24$ & $0\mathrm{R}25$ & 94 & $0\mathrm{R}183$ & $0\mathrm{L}208$ & 147 & $0\mathrm{R}159$ & $1\mathrm{R}12$ & 200 & $0\mathrm{L}195$ & $1\mathrm{R}195$\\
42 & $0\mathrm{R}13$ & $0\mathrm{L}16$ & 95 & $1\mathrm{L}109$ & $1\mathrm{L}95$ & 148 & $0\mathrm{R}83$ & $1\mathrm{R}143$ & 201 & $0\mathrm{R}202$ & $1\mathrm{R}198$\\
43 & $0\mathrm{R}55$ & $1\mathrm{R}56$ & 96 & $0\mathrm{R}19$ & $1\mathrm{R}110$ & 149 & $0\mathrm{R}144$ & $1\mathrm{R}155$ & 202 & $0\mathrm{L}123$ & $1\mathrm{R}168$\\
44 & $0\mathrm{R}57$ & $1\mathrm{R}58$ & 97 & $0\mathrm{R}60$ & $1\mathrm{R}60$ & 150 & $0\mathrm{L}66$ & $1\mathrm{L}141$ & 203 & $0\mathrm{R}197$ & $1\mathrm{R}28$\\
45 & $0\mathrm{R}24$ & $0\mathrm{R}54$ & 98 & $0\mathrm{R}68$ & $1\mathrm{R}135$ & 151 & $0\mathrm{R}150$ & $1\mathrm{R}170$ & 204 & $0\mathrm{R}163$ & $1\mathrm{R}193$\\
46 & $0\mathrm{R}60$ & $1\mathrm{R}61$ & 99 & $0\mathrm{R}93$ & $1\mathrm{R}108$ & 152 & $0\mathrm{R}159$ & $1\mathrm{R}114$ & 205 & $0\mathrm{R}146$ & $1\mathrm{R}177$\\
47 & $0\mathrm{R}62$ & $0\mathrm{L}16$ & 100 & $0\mathrm{R}105$ & $0\mathrm{L}208$ & 153 & $0\mathrm{R}97$ & $1\mathrm{R}97$ & 206 & $0\mathrm{R}161$ & $1\mathrm{R}194$\\
48 & $0\mathrm{R}25$ & $1\mathrm{R}48$ & 101 & $0\mathrm{R}117$ & $1\mathrm{R}118$ & 154 & $0\mathrm{R}49$ & $0\mathrm{R}32$ & 207 & $0\mathrm{R}45$ & $1\mathrm{R}100$\\
49 & $0\mathrm{L}35$ & $0\mathrm{R}63$ & 102 & $0\mathrm{R}71$ & $1\mathrm{R}58$ & 155 & $0\mathrm{R}192$ & $1\mathrm{R}188$ & 208 & $0\mathrm{L}4$ & $0\mathrm{L}208$\\
50 & $0\mathrm{R}16$ & $1\mathrm{R}50$ & 103 & $0\mathrm{R}57$ & $1\mathrm{R}119$ & 156 & $0\mathrm{R}150$ & $1\mathrm{R}112$ &  &  & \\
51 & $0\mathrm{R}147$ & $1\mathrm{R}143$ & 104 & $0\mathrm{R}21$ & $1\mathrm{R}120$ & 157 & $0\mathrm{L}187$ & $1\mathrm{R}96$ &  &  & \\
52 & $0\mathrm{R}67$ & $1\mathrm{R}68$ & 105 & $0\mathrm{R}21$ & $0\mathrm{L}10$ & 158 & $0\mathrm{R}168$ & $1\mathrm{R}179$ &  &  & \\
\bottomrule
\end{longtable}
\endgroup


\section{Provenance and acknowledgments}
The original U211 research packets and the unary-power-cover manuscript are
credited to the supplied ProveIt research collection. Their original notices,
nested dependency packets, and licenses are preserved unchanged in the
retained archives. The runtime engine and source-language utilities are
reused with their file-level provenance, not represented as newly authored
proof infrastructure.

The construction inherits the proof-search architecture associated with the
m254 reference development~\cite{m254}, the K215 verification
lineage~\cite{k215source,k215pa}, and the Theory K work~\cite{k230}.
Their respective fixed-machine theorems retain their original scope.
New work in this package is the comparison, compatible U210 composition,
shared-helper eight-factor reassociation, combined U209 fold, associated
proofs, independent audits, and reproducible final artifact assembly.

\begin{thebibliography}{9}
\bibitem{snapshot} V. Reshetnikov, \emph{ProveIt incoming research collection},
snapshot\\
{\small\hash{da4ab9dd68ccfc47ac852118601eaf1fdae53795}}.
\url{https://github.com/VladimirReshetnikov/ProveIt/tree/da4ab9dd68ccfc47ac852118601eaf1fdae53795/docs/incoming}.
The three exact U211 ZIP filenames and full hashes are in the package provenance.

\bibitem{unary} \emph{Unary power-cover theory and arithmetic calibration},
\texttt{pa\_unary\_power\_cover.zip}, in the same snapshot.
\url{https://github.com/VladimirReshetnikov/ProveIt/blob/da4ab9dd68ccfc47ac852118601eaf1fdae53795/docs/incoming/pa_unary_power_cover.zip}.
The original archive is included unchanged.

\bibitem{k215pa} E. Glazer, \emph{K215 halts iff ZFC is inconsistent
(Metamath Zero, PA)}. \url{https://github.com/elliotglazer/k215-pa}.

\bibitem{k215source} V. Reshetnikov, \emph{ProveIt K215 verification sources}.
\url{https://github.com/VladimirReshetnikov/ProveIt/tree/main/Computability/BusyBeaver/Research/ZFC/k225-formal/verification/k215}.

\bibitem{k230} E. Glazer, \emph{K230 and Theory K}.
\url{https://github.com/elliotglazer/k230-rocq/blob/main/README.md}.

\bibitem{m254} \emph{Busy-beaver ZF/m254 reference development}.
\url{https://github.com/jeff87654/busy-beaver/tree/zf-254/ZF/m254}.
\end{thebibliography}
\end{document}
